% Part: lambda-calculus
% Chapter: representability
% Section: currying

\documentclass[../../../include/open-logic-section]{subfiles}

\begin{document}

\olfileid{lam}{rep}{cur}
\olsection{کریئنگ}

اک $\lambd$-تجرید $\lambd[x][M]$ اک دلیل والا فنکشن ظاہر کردی اے۔ جدوں
اسیں کئی دلیلاں والا فنکشن متعین کرنا چاہیئے، ایہ حد بن جاندی اے۔اک طریقہ ایہ ہو سکدا سی کہ $\lambd$-کیلکولس
نوں جوڑے،تِن-ٹپل وغیرہ بناؤن دی اجازت دے کے ودھایا جائے؛فیر،
مثلاً،تِن رُکنی فنکشن $\lambd[x][M]$ اپنی دلیل وچ اک
تِن-ٹپل دی توقع رکھے گا۔پر ایہ کم \emph{کریئنگ} نال کرنا
زیادہ سہل اے۔

اک مثال ویکھئیے۔اک پل لئی فرض کر لیندے آں کہ $+$ دا عمل
$\lambd$-کیلکولس وچ ساڈے کول اے۔جمع دا فنکشن $2$-رُکنی اے،
یعنی دو دلیلاں لیندا اے۔پر $\lambd$-تجرید سانوں صرف اک
دلیل والے فنکشن دیندی اے:ایس دی نحو
$\lambd[(x,y)][(x+y)]$ ورگے اظہاراں دی اجازت نہیں دیندی۔اس
دے باوجود اسی اک دلیل والا فنکشن~$f_x(y)$ ویکھ سکدے آں جو
$\lambd[y][(x+y)]$ نال دتا گیا اے؛ایہ $x$ نوں اپنی اکلوتی
دلیل~$y$ وچ جوڑدا اے۔اصل وچ ایہ اک فنکشن نہیں،سگوں ``$x$ جوڑو''
والے وکھ وکھ فنکشناں دا خاندان اے،ہر عدد~$x$ لئی اک۔ہن اسی
اک ہور اک دلیل والا فنکشن~$g$ نوں $\lambd[x][f_x]$ دے طور اُتے
متعین کر سکدے آں۔دلیل $x$ اُتے لگن نال $g(x)$ فنکشن~$f_x$
واپس کردا اے—ایس لئی اوس دیاں قدراں آپ فنکشن نیں۔ہن جے
$g$ نوں $x$ اُتے لاگو کر کے حاصل نوں فیر~$y$ اُتے لاگو کریے،
تاں $(g(x))y = f_x(y) = x+y$ ملدا اے۔اینج اک دلیل والا
فنکشن~$g$ دو دلیل والے جمع فنکشن والا کم کر سکدا اے۔
``کریئنگ'' بس دو رُکنی فنکشن نوں اک دلیل والے فنکشوناں وچ بدلّن دی ترکیب دا ناں اے،
جنہاں دیاں قدراں آپ اک دلیل والے فنکشن ہوندے نیں۔

ہُن $\lambd$-کیلکولس دی اصل نحو وچ اک مثال ویکھئیے۔فنکشن
$f(x,y) = x$ نوں کیویں ظاہر کریے؟جے دو دلیلاں قبول کر کے پہلی
واپس کرن والا فنکشن تعریف کرنا ہووے،تاں
$\lambd[x][\lambd[y][x]]$ لکھ سکدے آں۔لفظی طور اُتے ایہ
فنکشن اک دلیل~$x$ قبول کر کے فنکشن~$\lambd[y][x]$ واپس کردا
اے۔فنکشن $\lambd[y][x]$ ہور دلیل~$y$ قبول کردا اے،پر اوس نوں
چھڈ دیندا اے تے ہمیشہ~$x$ واپس کردا اے۔ویکھئیے جدوں اسی
$\lambd[x][\lambd[y][x]]$ نوں دو دلیلاں اُتے لاگو کردے آں
تاں کی ہوندا اے:
\begin{align*}
  (\lambd[x][\lambd[y][x]])MN 
  \bredone & (\lambd[y][M])N \\
  \bredone & M
\end{align*}

عام طور اُتے،$x_1$، \dots،~$x_n$ پیرامیٹراں والے،کسی
ٹرم~$N$ نال تعریف کیتے فنکشن نوں اسی
$\lambd[x_1][\lambd[x_2][\ldots\lambd[x_n][N]]]$ لکھ سکدے
آں۔جے اوس نوں $n$ دلیلاں اُتے لاگو کریے تاں ایہ ملدا اے:
\begin{multline*}
  (\lambd[x_1][\lambd[x_2][\ldots\lambd[x_n][N]]]) M_1 \dots M_n \bredone\\
  \begin{aligned}
  \bredone {} & (\Subst{(\lambd[x_2][\ldots\lambd[x_n][N]])}{M_1}{x_1}) M_2
  \dots M_n\\
   \eqs {} & (\lambd[x_2][\ldots\lambd[x_n][\Subst{N}{M_1}{x_1}]]) M_2
            \dots M_n \\
  \vdots & \\
  \bredone {} &\Subst{\Subst{N}{M_1}{x_1}\ldots}{M_n}{x_n}
  \end{aligned}
\end{multline*}
آخری سطر دا لفظی مطلب ایہ اے کہ فنکشن دی تعریف دے اندرلے ٹرم وچ
$M_i$ نوں $x_i$ دی تھاں رکھیا جائے؛کئی دلیلاں والا فنکشن
لاگو کرن ویلے اسی عین ایہو چاہندے آں۔
\end{document}
